\pagenumbering{arabic}

\chapter{Introdução}
\label{ch:introducao}

% Meta editorial: 6--9 páginas. Este capítulo deve terminar com uma cadeia
% inequívoca entre problema, pergunta, hipótese, objetivo e validação.

\section{Contextualização}
% Introduzir progressivamente aprendizado contínuo, fluxos não estacionários,
% Class-Incremental Learning (CIL), cenário agnóstico à tarefa, esquecimento
% catastrófico e dilema estabilidade--plasticidade. Delimitar o trabalho à
% classificação de imagens. Não antecipar a formulação da ANT nesta seção.

\section{Motivação e justificativa}
% Partir das limitações práticas de CIL e da interferência durante atualizações
% incrementais. Apresentar o papel dos objetivos contrastivos e a possibilidade
% de continuarem exercendo pressão sobre relações já suficientemente
% preservadas. Preparar o conceito de ajustes não essenciais sem tratá-lo como
% causa de esquecimento já comprovada.

\section{Problema de pesquisa}
% Formular independentemente da solução: em métodos incrementais com
% destilação contrastiva, relações já suficientemente preservadas podem
% continuar contribuindo para o gradiente, produzindo atualizações
% potencialmente desnecessárias e interferência adicional ao aprender novas
% classes. Definir depois, no Capítulo 3, o que "suficientemente preservada" e
% "não essencial" significam operacionalmente.

\section{Questões de pesquisa}

\subsection{Questão principal}
% Investigar se restringir atualizações contrastivas às relações que ainda
% violam um critério de preservação reduz interferência sem prejudicar o
% aprendizado de novas classes e o desempenho incremental.

\subsection{Questões secundárias}
% Manter apenas questões respondidas pelos experimentos:
% (Q1) referência compartilhada ou por âncora na ANT;
% (Q2) cobertura intra-view ou simétrica completa;
% (Q3) normalização nGlobal/nLocal e sua equivalência efetiva;
% (Q4) comportamento das métricas internas do mecanismo;
% (Q5) efeito sobre accuracy e forgetting entre datasets e sementes.

\section{Hipóteses de pesquisa}

\subsection{Hipótese principal}
% H1: restringir a otimização contrastiva às violações relevantes reduz
% atualizações não essenciais sem comprometer a aquisição de novas classes,
% preservando ou melhorando o desempenho incremental. Formular também H0.

\subsection{Hipóteses auxiliares e critérios de interpretação}
% Criar uma hipótese auxiliar somente para decisões com comparação controlada.
% Para cada hipótese, registrar: métrica -> experimento -> evidência esperada ->
% critério para suportada, não suportada ou inconclusiva. Não considerar uma
% diferença pequena de uma única semente como confirmação isolada.

\section{Objetivos}

\subsection{Objetivo geral}
% Desenvolver e avaliar um método para evitar atualizações contrastivas não
% essenciais em CIL agnóstico à tarefa, tendo a ANT como proposta central.

\subsection{Objetivos específicos}
% Caracterizar o problema; formular a ANT; investigar referência e cobertura;
% auditar nGlobal/nLocal; desenvolver diagnósticos do mecanismo; avaliar em
% diferentes datasets/protocolos/sementes; comparar com TagFex e baselines
% adequados. Avg-K e SBS não são objetivos centrais.

\section{Escopo e delimitações}
% Fixar protocolo, informação disponível no treino/teste, memória limitada,
% backbones, datasets e métricas. Distinguir task-agnostic na inferência de
% task-free/boundary-free no treinamento. Separar trabalho consolidado de
% extensões ainda exploratórias.

\section{Contribuições esperadas}
% Hierarquia obrigatória:
% (1) contribuição principal: ANT;
% (2) decisões metodológicas: aGlobal/aLocal, intra-view/aSymFull e auditoria
%     nGlobal/nLocal;
% (3) formulações alternativas como ablações;
% (4) investigações auxiliares: Avg-K e SBS;
% (5) extensões: colisões entre tarefas e negativos baseados em protótipos.

\section{Organização do trabalho}
% Cap. 2: fundamentos, TagFex e lacuna; Cap. 3: ANT, decisões, validação e
% evidências preliminares; Cap. 4: trabalho restante até a defesa.
